Mesh-to-CAD reconstructors are scored on the geometry their search produces, but what they deliver is a program, and that program can fail to build or build an invalid solid. We evaluate systems by executing the delivered program in a sandbox, requiring every solid to be OCC-valid, and scoring the executed solid against the target, with any failed delivery scored zero. In three prospectively specified studies on hard-tier CADBench parts, this changes what works. On 300 held-out parts, a beam-search configuration of the CADFit optimizer obtains no valid replay on \the\numexpr300-\fDZeroValid\relax{} parts; delivering the best candidate program that executes to a valid solid raises validity-gated IoU by \hDXvsDZero{} (95\% interval \hDXvsDZeroCI; preregistered primary contrast), and three further stages add \hVTwovsDX, \hVThreevsVTwo{} and \hTrimvsVThree. Under the same judge, Claude Opus~5.5 writing CadQuery through a harness that executes every change scores \cOpusCode{} on 40 reserved held-out parts, \cOpusCodeVsVThree{} above the optimizer's v3 (exploratory). Offering the optimizer's stages to the agent as tools did not raise its score and produced deadline failures; a revised tool package with output-preserving rewrites and a per-check time cap, evaluated on 40 new parts with three episodes per condition, halves Opus's time with these tools (\tOpusSkillsVTwoMin{} to \tOpusSkillsVThreeMin{} minutes per part), every Opus delivery is valid, and the tools then beat plain code by \tOpusSkillsVsCode{} (Holm $p=\tOpusSkillsVsCodeHolm$). The package's preregistered quality gain over the original tools, \tOpusSkillsVersion, is not significant by its declared sign test.
